\chapter{Difference-in-Differences and Event Studies}
\label{ch:did-event-study}
\label{chap:did}

\begin{description}
  \item[Motivation.] Policies are often adopted at different times in different
  places.  Can these timing differences reveal causal effects?
  \item[Canonical DID.] Identify an ATT with two groups, two periods, and
  parallel trends.
  \item[Staggered adoption.] Define group--time effects when cohorts adopt in
  different periods.
  \item[Modern event studies.] Avoid already-treated controls, estimate
  transparent effects first, and aggregate second.
  \item[Credibility and inference.] Understand what pre-trends, confidence
  bands, and clustered standard errors can and cannot establish.
\end{description}

\section{Motivation}

Many important treatments in economics are policies rather than randomized
experiments.  A housing subsidy, minimum-wage increase, or new technology may
be introduced in one place before another.  A simple comparison between places
with and without the policy is unlikely to be causal: adopters may have had
different outcome levels before adoption.  A before--after comparison within
adopters is also insufficient: inflation, interest rates, or a common boom may
have changed outcomes even without the policy.

Difference-in-differences (DID) combines both comparisons:
\begin{enumerate}
  \item calculate the change over time for the treated group;
  \item calculate the change for an untreated comparison group; and
  \item subtract the second change from the first.
\end{enumerate}
Permanent level differences disappear in the first difference.  Common time
changes disappear in the second.  The remaining difference is causal under a
parallel-trends assumption.

\begin{example}[A two-by-two DID]
Suppose average monthly rent is
\begin{center}
\begin{tabular}{lrrr}
  \toprule
  Group & Before & After & Change\\
  \midrule
  Cities adopting a subsidy & 120 & 134 & 14\\
  Comparison cities         & 100 & 106 &  6\\
  \bottomrule
\end{tabular}
\end{center}
The DID is
\[
  (134-120)-(106-100)=14-6=8.
\]
Under parallel trends, 6 estimates how much rent would have risen in the
adopting cities without the subsidy.  The remaining 8 is the estimated effect
on those cities.
\end{example}

\section{Two groups and two periods}
\label{sec:did-two-by-two}

Let $t\in\{0,1\}$ and define
\begin{itemize}
  \item $G_i=1$ if unit $i$ belongs to the group treated beginning in period 1;
  \item $G_i=0$ if it belongs to the never-treated control group; and
  \item $D_{it}=G_i\1\{t=1\}$ as realized treatment status.
\end{itemize}
Let $Y_{it}(1)$ denote the potential outcome under the regime ``treated
beginning in period 1'' and let $Y_{it}(0)$ denote the outcome under ``never
treated.''  The observed regime is selected by group:
\begin{equation}
  Y_{it}=G_iY_{it}(1)+(1-G_i)Y_{it}(0).
  \label{eq:did-observed-simple}
\end{equation}
The target is the post-treatment average effect on the treated,
\begin{equation}
  ATT=\E[Y_{i1}(1)-Y_{i1}(0)\mid G_i=1].
  \label{eq:did-att-simple}
\end{equation}
DID most directly learns about the group that actually adopts.  The effect for
a never-adopting group may differ and is not identified by the same comparison.

\subsection{No anticipation}

Before the policy begins, the future treatment regime should not yet affect the
outcome:
\begin{equation}
  Y_{i0}(1)=Y_{i0}(0).
  \label{eq:did-no-anticipation-simple}
\end{equation}
This can fail if firms hire in advance of an announced subsidy or households
move before a program formally starts.  The last nominally pre-treatment period
may then already be treated in an economic sense.

\subsection{Parallel trends}

The key identifying assumption is
\begin{equation}
  \E[Y_{i1}(0)-Y_{i0}(0)\mid G_i=1]
  =\E[Y_{i1}(0)-Y_{i0}(0)\mid G_i=0].
  \label{eq:did-parallel-simple}
\end{equation}
This assumption concerns \emph{changes}, not levels.  Treated and control
groups may have different average outcomes before treatment, and selection on
permanent group differences is allowed.  The assumption concerns the untreated
potential outcome; after treatment, the treated group's counterfactual path is
unobserved and cannot be checked directly.

\subsection{Identification}

Under no anticipation and parallel trends,
\begin{align}
  ATT
  &=\E[Y_{i1}-Y_{i0}\mid G_i=1]
    -\E[Y_{i1}(0)-Y_{i0}(0)\mid G_i=1]\notag\\
  &=\E[Y_{i1}-Y_{i0}\mid G_i=1]
    -\E[Y_{i1}-Y_{i0}\mid G_i=0].
  \label{eq:did-identification-simple}
\end{align}
Every term in the final line is observed.  The sample estimator is
\begin{equation}
  \widehat{ATT}_{DID}
  =(\bar Y_{11}-\bar Y_{10})-(\bar Y_{01}-\bar Y_{00}),
  \label{eq:did-estimator-simple}
\end{equation}
where $\bar Y_{gt}$ is the sample mean for group $g$ in period $t$.  The order
of differencing is immaterial:
\[
  (\bar Y_{11}-\bar Y_{10})-(\bar Y_{01}-\bar Y_{00})
  =(\bar Y_{11}-\bar Y_{01})-(\bar Y_{10}-\bar Y_{00}).
\]

\subsection{Regression representation}

With two groups and two periods, Equation~\eqref{eq:did-estimator-simple} is the
OLS coefficient $\widehat\tau$ in
\begin{equation}
  Y_{it}=\beta_0+\beta_1G_i+\beta_2Post_t
         +\tau(G_i\times Post_t)+\varepsilon_{it},
  \label{eq:did-interaction-regression}
\end{equation}
where $Post_t=\1\{t=1\}$.  With panel data, the same coefficient comes from
\begin{equation}
  Y_{it}=\alpha_i+\lambda_t+\tau D_{it}+\varepsilon_{it},
  \label{eq:did-twfe-simple}
\end{equation}
where $\alpha_i$ and $\lambda_t$ are unit and time fixed effects.  In this
canonical setting, two-way fixed effects (TWFE) is simply another way to
calculate DID.

The same units are followed over time in panel data, whereas repeated cross
sections sample different units in each period.  Both can identify ATT under
appropriate sampling and parallel-trends assumptions.  Panel data are often
more efficient because differencing removes permanent unit-specific noise.

\section{Staggered treatment adoption}
\label{sec:did-staggered}

Now let $t=1,\ldots,T$.  Define $G_i=g$ if unit $i$ is first treated in period
$g$, and $G_i=\infty$ if it is never treated.  In R code, never-treated units
are usually recorded as $G_i=0$.  With absorbing treatment,
\begin{equation}
  D_{it}=\1\{t\geq G_i\}\quad\text{for eventually treated units},
  \qquad D_{it}=0\quad\text{if }G_i=\infty.
  \label{eq:did-staggered-status}
\end{equation}
Thus $D_{it}=1$ implies $D_{i,t+1}=1$.  A treatment that repeatedly switches
on and off requires a different design or additional assumptions.

Let $Y_{it}(g)$ be the outcome if unit $i$ is first treated in period $g$, and
let $Y_{it}(0)$ be the never-treated outcome.  No anticipation requires
\begin{equation}
  Y_{it}(g)=Y_{it}(0),\qquad t<g.
  \label{eq:did-no-anticipation-staggered}
\end{equation}
The natural causal building block is the group--time effect
\begin{equation}
  ATT(g,t)=\E[Y_{it}(g)-Y_{it}(0)\mid G_i=g],
  \qquad t\geq g.
  \label{eq:did-att-gt}
\end{equation}
$ATT(4,4)$ is the adoption-period effect for cohort 4; $ATT(4,6)$ is the effect
two periods later for the same cohort; and $ATT(6,6)$ concerns another cohort
in the same calendar period.  Policy effects can vary across all these cells.

\subsection{Valid control groups}

For $ATT(g,t)$, controls must be untreated in both the baseline $g-1$ and the
comparison period $t$.  Two choices are common:
\begin{description}
  \item[Never-treated controls.] Compare cohort $g$ with units that never adopt.
  \item[Not-yet-treated controls.] Compare cohort $g$ with units treated after
  $t$, together with never-treated units.
\end{description}
The second choice uses more observations but assumes later adopters provide a
valid counterfactual trend for earlier adopters.  This is an identification
choice, not merely a software option.

With never-treated controls, assume
\begin{equation}
  \E[Y_{it}(0)-Y_{i,g-1}(0)\mid G_i=g]
  =\E[Y_{it}(0)-Y_{i,g-1}(0)\mid G_i=\infty].
  \label{eq:did-pt-never}
\end{equation}
Then
\begin{equation}
  ATT(g,t)=
  \E[Y_{it}-Y_{i,g-1}\mid G_i=g]
  -\E[Y_{it}-Y_{i,g-1}\mid G_i=\infty].
  \label{eq:did-id-never}
\end{equation}
For not-yet-treated controls, replace $G_i=\infty$ by $G_i>t$ in both displays,
using the convention that $\infty>t$.  Each $ATT(g,t)$ is then a clean
two-group, two-period DID.  Already-treated observations are not used as
untreated controls.

\subsection{Conditional parallel trends and overlap}

Let $W_i$ contain pre-treatment covariates.  A conditional assumption for
not-yet-treated controls is
\begin{align}
  &\E[Y_{it}(0)-Y_{i,g-1}(0)\mid G_i=g,W_i]\notag\\
  &\qquad=\E[Y_{it}(0)-Y_{i,g-1}(0)\mid G_i>t,W_i].
  \label{eq:did-conditional-pt}
\end{align}
Overlap requires controls at every covariate value represented in cohort $g$:
\begin{equation}
  \Pr(G_i>t\mid W_i=w)>0
  \quad\text{for }w\in\operatorname{supp}(W_i\mid G_i=g).
  \label{eq:did-overlap}
\end{equation}
Under these conditions,
\begin{align}
 ATT(g,t)=\E\Big[&
   \E[Y_{it}-Y_{i,g-1}\mid G_i=g,W_i]
   \notag\\[-2pt]
 &-\E[Y_{it}-Y_{i,g-1}\mid G_i>t,W_i]
   \ \Bigm|\ G_i=g\Big].
 \label{eq:did-conditional-id}
\end{align}
Outcome regression estimates conditional control trends; propensity weighting
reweights controls toward cohort $g$; and a doubly robust DID combines both.
Under the relevant regularity conditions, a doubly robust estimator remains
consistent if either the outcome-trend model or the generalized propensity
model is correct.  Covariates should be measured before treatment, not caused
by it.

\section{Why conventional TWFE can fail}
\label{sec:did-twfe-failure}

With many periods, researchers often estimate
\begin{equation}
  Y_{it}=\alpha_i+\lambda_t+\beta D_{it}+u_{it}.
  \label{eq:did-static-twfe}
\end{equation}
Under staggered adoption and heterogeneous effects, $\beta$ need not be an
interpretable average treatment effect.

Imagine an early-treated group E, a late-treated group L, and a never-treated
group N.  When E is newly treated, L and N are untreated.  When L is newly
treated, E is already treated.  Conventional TWFE can use all of the following:
\begin{enumerate}
  \item newly treated versus never treated;
  \item newly treated versus not yet treated; and
  \item newly treated versus already treated.
\end{enumerate}
The third comparison subtracts a change containing E's treatment-effect
dynamics.  It is not L's untreated counterfactual trend.

The Goodman--Bacon decomposition gives nonnegative weights to component
two-by-two DID comparisons, but some components use already-treated controls.
When the TWFE coefficient is instead written in terms of cohort--time treatment
effects, implicit weights can be negative.  In extreme cases, every causal
effect can be positive while the static TWFE coefficient is negative.

\subsection{Conventional event-study regressions}

A familiar dynamic regression is
\begin{equation}
  Y_{it}=\alpha_i+\lambda_t+
  \sum_{e\neq-1}\mu_e\1\{t-G_i=e\}+u_{it},
  \label{eq:did-conventional-event-study}
\end{equation}
where $e=t-G_i$ is event time and $e=-1$ is omitted.  With staggered timing and
heterogeneous effects, $\mu_e$ can be contaminated by effects at other event
times.  Even a lead can be nonzero because of post-treatment heterogeneity
rather than a genuine pre-trend.  Adding leads and lags to TWFE does not by
itself repair the control-group problem.

Fixed effects remain useful.  TWFE equals DID in the canonical two-by-two
design; a single treated cohort with a valid never-treated group avoids
already-treated controls; and a homogeneous common-effect model can justify a
static coefficient.  The lesson is to define the causal parameter and valid
controls before running the regression.

\section{Modern DID and event studies}
\label{sec:did-modern}

\begin{important}
Estimate transparent cohort--time causal effects using untreated controls, and
only then aggregate them into the summary required by the empirical question.
\end{important}

The Callaway--Sant'Anna approach estimates the collection of $ATT(g,t)$ values
using clean comparisons.  Depending on the design and covariates, each cell can
be estimated by differences in outcome changes, outcome regression, inverse
probability weighting, or doubly robust DID.

\subsection{Aggregating group--time effects}

One overall summary is
\begin{equation}
  \theta^{overall}=\sum_g\sum_{t\geq g}
  \omega_{g,t}ATT(g,t),
  \qquad \omega_{g,t}\geq0,\quad
  \sum_{g,t}\omega_{g,t}=1.
  \label{eq:did-overall-aggregation}
\end{equation}
Other summaries average within adoption cohort or within calendar period.  They
answer different questions.  Software cannot decide which is economically
relevant.

In the \texttt{did} package, \texttt{type = "simple"} weights post-treatment
cells in proportion to cohort size, so early cohorts receive more total weight
because they contribute more cells.  \texttt{type = "group"} first averages
post-treatment effects within cohort.  \texttt{type = "dynamic"} averages
contributing cohorts at each event time, using cohort-size weights.

For event time $e=t-g$, define
\begin{equation}
  \theta(e)=\sum_{g:g+e\leq T}
  \Pr(G_i=g\mid G_i+e\leq T)ATT(g,g+e).
  \label{eq:did-dynamic-aggregation}
\end{equation}
For $e\geq0$, this is the average effect after $e$ periods of exposure among
cohorts observed at that horizon.  At long horizons only early adopters remain,
so a changing curve can reflect both genuine dynamics and changing cohort
composition.  A balanced event window holds the contributing cohorts fixed,
improving comparability at the cost of discarding data.

Other modern approaches implement a similar principle in different ways:
\begin{description}
  \item[Sun--Abraham.] Estimate cohort-by-event-time interactions and aggregate
  them without contamination from other event times.
  \item[Borusyak--Jaravel--Spiess.] Fit an untreated-outcome model using
  untreated observations, impute untreated counterfactuals for treated cells,
  and average the resulting effects.
  \item[Group--time DID.] Construct clean two-by-two comparisons for each
  $(g,t)$ and then aggregate.
\end{description}
These methods can target different parameters and maintain different outcome,
comparison-group, or covariate models.  They are not interchangeable commands.

\section{Credibility and inference}
\label{sec:did-inference}

Modern event studies often report placebo or pseudo-ATT estimates before
treatment.  Large, systematic pre-treatment differences weigh against the
proposed parallel-trends design.  Estimates near zero do not prove it:
\begin{itemize}
  \item pre-treatment estimates may be noisy and have low power;
  \item paths may be parallel before treatment but diverge afterward;
  \item choosing a specification because its pre-trend test is insignificant
        changes the statistical behavior of the final estimates; and
  \item the assumption concerns an unobserved post-treatment counterfactual.
\end{itemize}
Pre-trends are a diagnostic, not a certificate of validity.

An event study reports many coefficients.  Pointwise 95\% intervals cover each
coefficient separately, not the complete curve jointly.  For statements such
as ``all pre-treatment effects are zero,'' simultaneous confidence bands are
preferable.  Modern DID packages can obtain uniform bands by multiplier
bootstrap.

Outcomes from the same unit are dependent over time.  When treatment is assigned
at a higher level, observations may also be dependent within that assignment
cluster.  With state-level policies, cluster at the state level rather than the
individual level.  With few clusters, ordinary cluster-robust approximations
may be unreliable.  The clustering level follows treatment assignment and the
dependence structure, not the number of rows.

\subsection{A design checklist}

Before estimating a staggered DID, ask:
\begin{enumerate}
  \item What is the unit, and when is it first treated?
  \item Is treatment absorbing, or can it turn off?
  \item Could units anticipate treatment?
  \item Which observations are valid controls for each cohort and period?
  \item Is parallel trends more plausible conditionally on pre-treatment
        covariates?
  \item Is the target group--time, dynamic, calendar-time, or overall ATT?
  \item Does cohort composition change across the event window?
  \item At what level should inference be clustered?
  \item Do other policies or shocks occur at the same time?
  \item Are conclusions robust to outcome scale, sample window, control group,
        and anticipation window?
\end{enumerate}

\section{R implementation}
\label{sec:did-r}

We simulate a balanced panel in which cohorts are first treated in periods 4,
6, or 8, while some units are never treated.  Treatment effects grow with
exposure and differ across cohorts---exactly the setting in which conventional
TWFE can be hard to interpret.

\subsection{Generate the panel}

\begin{rcode}
set.seed(5280)
N <- 600
TT <- 10

units <- data.frame(
  id = 1:N,
  G = sample(c(0, 4, 6, 8), N, replace = TRUE,
             prob = c(0.30, 0.25, 0.25, 0.20)),
  alpha = rnorm(N),
  W = rnorm(N)
)

data <- merge(expand.grid(id = 1:N, t = 1:TT), units,
              by = "id", sort = TRUE)
data$D <- as.integer(data$G > 0 & data$t >= data$G)
data$event_time <- ifelse(data$G > 0, data$t - data$G, NA)

data$cohort_scale <- ifelse(data$G == 4, 1.50,
                     ifelse(data$G == 6, 1.00,
                     ifelse(data$G == 8, 0.60, 0)))
data$tau <- ifelse(data$D == 1,
                   data$cohort_scale * (1 + 0.5 * data$event_time), 0)

data$lambda <- 0.25 * data$t + 0.03 * data$t^2
data$Y0 <- data$alpha + data$lambda + 0.5 * data$W +
           rnorm(nrow(data))
data$Y <- data$Y0 + data$tau
\end{rcode}

Untreated outcomes satisfy parallel trends by construction.  The treatment
effect is positive but heterogeneous across cohorts and exposure lengths.

\subsection{Calculate one group--time DID by hand}

\begin{rcode}
simple_att_gt <- function(data, g, t) {
  stopifnot(t >= g)
  base <- g - 1

  now <- data[data$t == t, c("id", "G", "Y")]
  names(now)[names(now) == "Y"] <- "Y_now"
  before <- data[data$t == base, c("id", "Y")]
  names(before)[names(before) == "Y"] <- "Y_before"

  change <- merge(now, before, by = "id")
  change$dY <- change$Y_now - change$Y_before
  treated <- change$G == g
  controls <- change$G == 0 | change$G > t

  mean(change$dY[treated]) - mean(change$dY[controls])
}

simple_att_gt(data, g = 4, t = 6)
mean(data$tau[data$G == 4 & data$t == 6])
\end{rcode}

The two numbers should be close.  The second is available only because the DGP
is known; in real data it is the unobserved causal parameter.

\subsection{Callaway--Sant'Anna group--time effects}

\begin{rcode}
library(did)

att <- att_gt(
  yname = "Y", tname = "t", idname = "id", gname = "G",
  xformla = ~ 1, data = data, panel = TRUE,
  control_group = "notyettreated", est_method = "dr",
  base_period = "universal", clustervars = "id",
  bstrap = TRUE, cband = TRUE,
  biters = 199                 # short browser demonstration only
)

summary(att)
ggdid(att)
\end{rcode}

With \texttt{xformla = \textasciitilde{} 1}, the built-in outcome-regression,
IPW, and doubly robust implementations have the same unconditional two-by-two
DID point estimates.  Replace it by, for example,
\texttt{\textasciitilde{} W} for a conditional design.

The universal base period uses $g-1$ throughout and normalizes event time $-1$
to zero.  The varying-base default leaves post-treatment effects unchanged but
uses adjacent-period comparisons for pre-treatment pseudo-ATTs.  The multiplier
bootstrap creates simultaneous bands.  We use only 199 repetitions to keep a
live classroom run short; use at least 999, and often more, for final empirical
inference.  If treatment is assigned above the unit level, specify the relevant
assignment cluster.  For panel data, \texttt{did} permits at most two cluster
variables and requires one to be the unit identifier.

\subsection{Aggregate the effects}

\begin{rcode}
overall <- aggte(att, type = "simple")
by_group <- aggte(att, type = "group")
dynamic <- aggte(att, type = "dynamic", min_e = -3, max_e = 4)

summary(overall)
summary(by_group)
summary(dynamic)
ggdid(dynamic)

dynamic_balanced <- aggte(
  att, type = "dynamic",
  min_e = -3, max_e = 3,
  balance_e = 3
)
summary(dynamic_balanced)
ggdid(dynamic_balanced)
\end{rcode}

The \texttt{simple}, \texttt{group}, and \texttt{dynamic} outputs answer
different questions and need not be numerically equal.

\subsection{Conventional and interaction-weighted event studies}

\begin{rcode}
library(fixest)

data$ever_treated <- as.integer(data$G > 0)
data$event_twfe <- ifelse(data$G > 0, data$t - data$G, 0)

twfe_event <- feols(
  Y ~ i(event_twfe, ever_treated, ref = -1) | id + t,
  data = data, cluster = ~ id
)

# A cohort outside the observed period range is treated as never treated.
data$G_sunab <- ifelse(data$G == 0, 1000, data$G)
sa_event <- feols(
  Y ~ sunab(G_sunab, t, ref.p = -1) | id + t,
  data = data, cluster = ~ id
)

iplot(list("Conventional TWFE" = twfe_event,
           "Sun--Abraham" = sa_event),
      ref.line = -1, main = "Event-study estimates")
\end{rcode}

The \texttt{fixest::iplot()} intervals here are pointwise.  They are not the
multiplier-bootstrap simultaneous bands shown by \texttt{ggdid()}.

The companion page pins WebR 0.5.9 (R 4.5), for which the WebAssembly builds of
\texttt{did}, its \texttt{fastglm} dependency, and \texttt{fixest} are mutually
compatible.  This is a runtime compatibility pin, not an econometric choice.

\section*{Summary}

\begin{itemize}
  \item DID compares outcome changes, not outcome levels.
  \item No anticipation and parallel trends for untreated potential outcomes
        identify the canonical two-by-two ATT.
  \item Under staggered adoption, $ATT(g,t)$ is the transparent causal building
        block; never-treated or not-yet-treated units can be controls.
  \item Conventional static and dynamic TWFE can be contaminated by
        already-treated comparisons when effects are heterogeneous.
  \item Modern procedures estimate valid cohort--time comparisons first and
        aggregate them second.
  \item Pre-trends, aggregation, cohort composition, anticipation, covariates,
        and clustering are substantive design choices.
\end{itemize}

\section*{Further reading}

\begin{itemize}
  \raggedright
  \item Callaway and Sant'Anna (2021), ``Difference-in-Differences with Multiple
  Time Periods,'' \emph{Journal of Econometrics},
  \url{https://doi.org/10.1016/j.jeconom.2020.12.001}.
  \item Sun and Abraham (2021), ``Estimating Dynamic Treatment Effects in Event
  Studies with Heterogeneous Treatment Effects,'' \emph{Journal of
  Econometrics}, \url{https://doi.org/10.1016/j.jeconom.2020.09.006}.
  \item Goodman-Bacon (2021), ``Difference-in-Differences with Variation in
  Treatment Timing,'' \emph{Journal of Econometrics},
  \url{https://doi.org/10.1016/j.jeconom.2021.03.014}.
  \item Borusyak, Jaravel, and Spiess (2024), ``Revisiting Event-Study Designs:
  Robust and Efficient Estimation,'' \emph{Review of Economic Studies},
  \url{https://doi.org/10.1093/restud/rdae007}.
  \item Roth, Sant'Anna, Bilinski, and Poe (2023), ``What's Trending in
  Difference-in-Differences?,'' \emph{Journal of Econometrics},
  \url{https://doi.org/10.1016/j.jeconom.2023.03.008}.
  \item Sant'Anna and Zhao (2020), ``Doubly Robust Difference-in-Differences
  Estimators,'' \emph{Journal of Econometrics},
  \url{https://doi.org/10.1016/j.jeconom.2020.06.003}.
  \item Rambachan and Roth (2023), ``A More Credible Approach to Parallel
  Trends,'' \emph{Review of Economic Studies},
  \url{https://doi.org/10.1093/restud/rdad018}.
\end{itemize}
